
on peut dans ce diagramme remplacer les @@M0@@ par les @@M1@@, et on trouve un diagramme du type voulu, en termes des anneaux de Chow de @@M2@@ et @@M3@@.

\textbf{Remarque 7.13.13.} Soit @@M4@@ un schéma lisse et projectif sur un corps algébriquement clos. Il résulte trivialement du formalisme développé dans (SGA 5 IV) que si @@M5@@ est un cycle de codimension @@M6@@ sur @@M7@@ dont l'image dans @@M8@@ est nulle, alors @@M9@@ est numériquement équivalent à zéro. C'est une des conjectures standard de la théorie des cycles algébriques que la réciproque est également vraie. Il est immédiat d'ailleurs grâce à (VII 4.11) que l'on trouve une conjecture équivalente en y remplaçant le groupe des cycles algébriques par le groupe @@M10@@, en utilisant la définition (7.13.8), ou encore par le groupe @@M11@@, en écrivant que l'équivalence numérique à zéro d'un élément de @@M12@@ est équivalente à la nullité de son augmentation et de ses classes de Chern dans les @@M13@@.

Ceci posé, il résulte immédiatement de la commutativité de (7.13.10) que si @@M14@@ satisfait à la conjecture précédente, alors l'homomorphisme de spécialisation @@M15@@ (et a fortiori @@M16@@) est compatible avec l'équivalence numérique, comme il a été conjecturé dans 7.9. On notera d'ailleurs que la conjecture à laquelle on a fait appel pour cela (équivalence numérique implique équivalence @@M17@@-cohomologique) est plus faible que celle envisagée dans 7.12.5.

\subsection*{7.14. Homomorphisme de spécialisation pour les groupes de cycles et classes de cycles.}
Si @@M18@@ est un schéma noethérien, désignons par @@M19@@ (resp. @@M20@@) le groupe des cycles algébriques sur @@M21@@ purement de dimension @@M22@@ (resp. purement de codimension @@M23@@). Lorsque @@M24@@ est de type fini, on définit alors
